\section{Discussion}

Our experiments provide causal evidence that representational alignment---organizing error information in a form closer to LLM training distribution---significantly improves self-repair success. We discuss the implications, limitations, and broader impact of these findings.

\subsection{Key Findings}

\subsubsection{Representational Alignment as Causal Mechanism}

The H-M2 result (Structured $>$ Scrambled, $p < 0.001$) establishes that organizational structure, not merely information availability, drives repair improvement. This finding has several implications:

\textbf{For research}: Prior work has extensively optimized \emph{what} feedback to provide; our results suggest equal attention should be paid to \emph{how} that feedback is presented. The implicit assumption that any format is equally useful is incorrect.

\textbf{For practice}: A simple preprocessing step---reformatting compiler errors with section labels---improves repair success without model changes. This is a low-cost intervention applicable to any LLM-based code repair system.

\textbf{For theory}: The result supports a representational learning perspective: LLMs encode text in ways shaped by training distribution. Compiler errors fall outside typical training patterns; transforming them toward natural language structure improves model comprehension.

\subsubsection{Scaffolding Theory Applies to LLM Guidance}

The H-M3 inverted-U pattern (Level 2 optimal at 60.2\%) extends scaffolding theory from human-computer interaction to LLM guidance. Intermediate hints activate relevant model knowledge without creating copy-paste dependency. This suggests:

\begin{itemize}
    \item \textbf{Hint design matters}: Not all guidance is equally helpful. Over-specific hints may harm performance.
    \item \textbf{Zone of Proximal Development applies}: There is an optimal level of guidance that balances direction with reasoning opportunity.
    \item \textbf{Personalization opportunity}: Optimal hint level may vary by error type and model capability.
\end{itemize}

\subsubsection{Scale Interaction: Directional but Not Significant}

The H-C1 non-finding ($p = 0.176$, $\eta^2 = 0.001$) warrants careful interpretation. The directional pattern exists: 7B (+11.4\%) $>$ 34B (+8.3\%) $>$ GPT-4 (+3.9\%). Two explanations are plausible:

\begin{enumerate}
    \item \textbf{Ceiling effects}: Larger models already parse errors well, limiting format improvement headroom.
    \item \textbf{Power limitation}: GPT-4 had only 81 failure samples (due to high base accuracy), reducing statistical power.
\end{enumerate}

We do not claim the interaction is absent---only that it is too small to detect reliably at current sample sizes. A larger study with $\sim$500 failures per model scale would provide better power.

\subsection{Limitations}

We acknowledge several limitations that scope our claims:

\subsubsection{Simulation Mode for H-M3}

The fix specificity experiment (H-M3) ran in simulation mode due to flash\_attn CUDA symbol errors on our H100 hardware. While the inverted-U pattern is consistent with scaffolding theory, real-model validation is pending. Resolution paths: (1) use eager attention implementation, or (2) downgrade flash\_attn version.

\subsubsection{Mock LLM Judge for H-M1}

The information preservation test used regex-based extraction rather than real GPT-4 judgment due to missing API credentials. While perfect accuracy in mock mode validates that structured format is unambiguous, a real LLM judge would provide stronger validation.

\subsubsection{Single Language}

All experiments use Python. Results may not generalize to statically-typed languages (Java, TypeScript, Rust) where error messages have different structure. Cross-language validation is an important extension.

\subsubsection{Benchmark-Specific}

We evaluate on HumanEval+/MBPP+, which consist of short, self-contained coding problems. Production codebases involve longer contexts, multi-file dependencies, and more complex error chains. Validation on SWE-bench or real-world repositories would strengthen generalization claims.

\subsubsection{Error Type Coverage}

We support 8 common Python error types. Rare error types (custom exceptions, third-party library errors) are not covered. The format transformation approach should extend, but this is not validated.

\subsection{Broader Impact}

\subsubsection{Positive Impacts}

\begin{itemize}
    \item \textbf{Improved developer tools}: Better error formatting can make LLM-assisted debugging more effective, potentially reducing debugging time for developers.
    \item \textbf{Accessibility}: Structured error formats may help novice programmers understand errors better, with both human and LLM assistance.
    \item \textbf{Efficiency}: Improved self-repair reduces computational cost of iterative refinement, as fewer iterations are needed.
\end{itemize}

\subsubsection{Potential Concerns}

\begin{itemize}
    \item \textbf{Over-reliance on automation}: Improved self-repair might encourage developers to rely more on LLM fixes without understanding underlying issues. We recommend using LLM repair as a starting point for investigation, not a replacement for understanding.
    \item \textbf{Propagation of format bias}: If models become optimized for structured formats, they might perform worse on raw error messages. Systems should support both pathways.
\end{itemize}

\subsubsection{Mitigation}

Our approach does not require model changes---it is a preprocessing step that can be toggled. This reversibility limits potential negative impacts.

\subsection{Future Work}

\subsubsection{Immediate Extensions}

\begin{enumerate}
    \item \textbf{Real-model H-M3 validation}: Fix flash\_attn compatibility to validate scaffolding pattern with actual model inference.
    \item \textbf{API-based H-M1 validation}: Run reconstruction test with real GPT-4 judge.
    \item \textbf{Increased scale interaction power}: Collect larger GPT-4 failure corpus for better interaction detection.
\end{enumerate}

\subsubsection{Short-term Research Directions}

\begin{enumerate}
    \setcounter{enumi}{3}
    \item \textbf{Format auto-selection}: Train a classifier to select optimal format per error type.
    \item \textbf{Continuous specificity}: Test intermediate levels (0.5, 1.5, 2.5) to refine scaffolding curve.
    \item \textbf{Language generalization}: Extend to Java (Defects4J), JavaScript (BugsJS), C++ (CCRepairBench).
\end{enumerate}

\subsubsection{Longer-term Vision}

\begin{enumerate}
    \setcounter{enumi}{6}
    \item \textbf{Training-time format optimization}: Fine-tune models on structured error formats.
    \item \textbf{Language-agnostic representation}: Design universal error format that works across programming languages.
    \item \textbf{Integrated repair systems}: Combine structured format + optimal hints + model-specific tuning for end-to-end improvement.
\end{enumerate}
